任务执行中，worker 可能完成计算却来不及报告，
也可能报告已到达而确认回复丢失。
调度者重新运行任务时，两次尝试可以同时存在。
输出协议把每次尝试产生的字节与逻辑任务的正式输出关联起来，
使计算重执行有明确的收束位置。

\subsection{临时文件与胜出尝试}

\code{M3-a1} 写入自己的临时目录，
完成后报告文件集合、长度和必要的校验信息。
master 若尚未选定 $M_3$ 的完成尝试，就记录这次输出。
若已经选择 \code{M3-a2}，就忽略迟到的 \code{a1} 完成报告。
reduce 读取被选中的集合。
一个 map 尝试产生多个 reduce 片段时，
master 一次记录这一尝试的整个片段集合，
使后续读取有明确的来源~\cite{mapreduce-design}。

推测执行主动制造这种并存。
某个任务明显慢于同阶段其他任务，
调度者可以在另一台机器上启动相同逻辑任务的备用尝试。
它们读取同一输入、执行同一计算，其中一个成为正式完成者。
慢尝试可能随后结束，甚至继续运行已经发出的代码。
取消命令用于释放执行资源，输出选择用于决定读者采用哪一份结果。

\CTwentyFigure{ch20-04}{推测执行允许多个尝试运行，输出入口选择逻辑任务的一份结果}{fig:ch20-speculation}

MapReduce 原设计对 reduce 使用底层文件系统的原子 \code{rename}，
把私有临时输出换成最终文件。
多个尝试可能竞争这个名字，
最终文件包含其中一次完整执行产生的字节。
该结论依赖所用文件系统的重命名语义及一个 reduce 的输出布局。
应用额外生成多个有相互一致性要求的文件时，
每次重命名只能解释相应文件的可见变化，
整组文件的可见入口要由确定性计算或另外的发布协议组织。

确定性计算使不同尝试的输出可以互换。
若计算读取当前价格、当前时间或未固定的外部服务结果，
两次尝试可能产生不同内容。
即使每个最终文件都完整来自一个尝试，
多个输出分区也可能选择彼此不兼容的观察。
需要共同一致结果的作业将变化输入固定到作业清单，
或者通过明确的事务和版本规则选择一套观察。

中间结果丢失后的重算又涉及位置变化。
调度者得知旧片段不可取后，安排新尝试产生同一逻辑任务的输出，
再更新下游读取的位置。
记录输入版本、任务代次和正式片段集合，
可以区分当前重算的完成报告与上一代的迟到报告。
这种显式代次是本书发布设计中的关联字段，
具体 MapReduce 产品如何拒绝旧尝试，则由它自己的任务与提交协议决定。

\subsection{Hadoop 的存储、资源与作业角色}

Hadoop 是面向集群的大数据软件体系。
其中 HDFS 保存分布式文件，YARN 分配集群执行资源，
Hadoop MapReduce 在这些资源上组织 map、shuffle 和 reduce。
第十九章已经说明 HDFS 的 NameNode 管理文件与块位置，
DataNode 保存块副本。
计算端利用这些位置选择靠近数据的机器，
失败任务则重新读取原输入，在新位置计算~\cite{hadoop342-mapreduce,hadoop342-yarn}。

YARN 将集群资源分配与每个应用的任务管理分开。
ResourceManager 在应用之间分配资源，
NodeManager 是每台机器上的代理，负责启动和监测本机执行单元。
ApplicationMaster 是每个应用的管理进程，
向 ResourceManager 申请资源，再与 NodeManager 协作启动和跟踪任务。
Hadoop MapReduce 使用自己的 MRAppMaster 承担这个应用级角色。
因此，Google MapReduce 原设计中的一个 master 角色，
在这里分成了集群资源管理与具体作业管理两个层次。

YARN 的 container 表示一次资源分配，
包括内存、CPU 等运行条件以及相应的执行上下文。
NodeManager 据此启动任务进程。
这里的资源容器并不要求采用 Docker 镜像。
一个任务需要重新执行时，ApplicationMaster 请求新的执行资源，
逻辑任务和固定输入保持原关联，机器、进程和尝试编号可以变化。
图~\ref{fig:ch20-hadoop-layout} 把存储、副本位置、资源申请与任务字节路径放在一起。

\CTwentyFigure{ch20-26}{Hadoop 将文件保存、集群资源分配和每作业任务管理分开}{fig:ch20-hadoop-layout}

\subsection{Hadoop 提交器中的两个层次}

任务进程写出文件后，提交器负责把尝试输出放到规定位置。

Hadoop 3.4.2 的 \code{FileOutputCommitter}
区分任务尝试目录、任务已提交目录与作业输出目录，
并提供两种算法~\cite{hadoop342-committer}。
算法 1 在任务提交时，将尝试目录移动到该任务的已提交位置，
到作业提交时再汇合各任务。
算法 2 在任务提交时，直接把文件合并到最终输出目录，
作业收尾主要进行清理和成功标记。
该版本支持两条路径，默认选择算法 2。

任务已经提交与整个作业已经形成完整结果集合，有各自的观察条件。
算法 2 中，作业尚未完成时，最终目录可能已经出现部分任务文件。
读取者若只枚举现有文件，就会把尚未完成的集合当成输入。
读取协议因此先确定要读哪个作业版本，再确认完成条件及应有的分区集合。

\CTwentyFigure{ch20-05}{Hadoop 两种提交路径把任务提交与作业完成放在不同位置}{fig:ch20-hadoop-commits}

\code{\_SUCCESS} 是作业完成协议中的标记，是否创建还受配置影响。
它帮助读者判断相应收尾动作是否完成。
需要长期审计的对账结果，还可以发布一份清单，
列明来源、输入摘要、输出分区及各文件摘要。
读取者从这份清单取得确定的结果集合，随后逐项检查文件，
便能发现缺分区、错误版本或内容损坏。

\code{abortTask} 清理尝试目录，减少废弃文件。
清理尚未引用的尝试文件与撤回已提交业务效果，是不同的动作。
文件系统的目录重命名与对象存储的名称变换也要分别处理。
有些对象接口中的 rename 由复制、再删除组成，
跨对象操作没有 HDFS 目录重命名相同的原子范围。
对象存储提交器需要按具体连接器组织上传、完成上传和发布，
读者依据该协议认可的入口取得完整集合。

本书的对账发布可以沿用第十九章的先数据、后清单结构。
每个尝试写入自己的不可变输出对象，
提交者验证分区齐全及摘要，再发布带 $J_1$ 身份的 manifest。
若发布入口支持条件创建，提交者可以在尚无该 $J_1$ 清单时绑定整组对象。
发布回复丢失后，按 $J_1$ 重读并比较完整清单。
已有内容相同就返回原发布结果，内容不同则报告冲突。
这些动作依赖发布入口实际提供的原子绑定与持久性，
底层对象字节的上传和废弃对象清理各自按存储协议进行。

\CTwentyFigure{ch20-27}{不可变输出先到位，作业清单再一次绑定完整结果集合}{fig:ch20-output-manifest}

\subsection{Spark 的数据集、依赖与执行进程}

Spark 是通用的分布式数据处理引擎，
允许应用把读取、过滤、映射、分组和汇总连接成多步计算。
同一批数据可以在多个后续步骤中复用，
无须把每一步都组织成独立的 map/reduce 作业。
它也能在 YARN 或其他集群管理器上取得资源，
计算引擎与资源管理器分别负责应用内部执行和应用之间的资源分配。
以下结合 Spark 4.1.3 的 RDD 与输出路径说明~\cite{spark413-rdd,spark413-cluster}。

RDD，即弹性分布式数据集，是按分区分布在节点上的数据集合。
应用对已有 RDD 施加变换，得到新的 RDD；
原集合与变换关系保存在计算依赖中。
例如，固定账本先经过解析与格式核对，
再按平台身份分组，最后形成差异结果。
仅建立这些变换时，Spark 可以先保存计算关系，
等保存结果或取回汇总等 action 真正要求结果时，再安排任务执行。

有些变换只依赖某个上游分区，记录可在同一任务内连续处理。
按付款身份重新分组则涉及多个上游分区：
同一键的记录分散在不同机器，需要先重新分区并交换字节。
这一步称为 shuffle，它在执行图中形成阶段边界。
stage 是一组可按当前分区依赖一起调度的任务，
每个 task 处理其中一个输出分区。
driver 是运行应用主程序并组织这些阶段的进程，
executor 是为该应用启动、在工作节点上计算和保存缓存数据的进程。
不同 task attempt 则表示同一任务的不同执行尝试。

\CTwentyFigure{ch20-28}{RDD 保存变换依赖，shuffle 汇合同键记录并分开执行阶段}{fig:ch20-spark-dag}

RDD 分区丢失后，可以沿依赖重新计算。
缓存使已算出的分区可直接复用，减少多次 action 的计算成本。
可靠检查点把选定数据集写入配置的可靠存储，
恢复可以从该数据集继续，减少必须重走的依赖链。
\code{localCheckpoint} 使用 executor 本地临时存储，
它在截断原依赖后若再丢失这份数据，作业可能无法恢复。
因此，对账输入与正式差异输出分别发布持久版本，
运行中间分区再按缓存、重算或检查点的条件恢复。

\subsection{Spark 的重试、推测与提交协调}

一个分区的不同 attempt 可以先后重试，
也可以因推测执行而并存。

Spark 4.1.3 的 \mbox{\code{spark.speculation}} 默认关闭，
开启后结合阶段完成比例、运行时长和效率等条件判断慢任务。
\mbox{\code{spark.task.maxFailures}} 默认 4，
表示特定任务连续失败达到相应次数后放弃，
对应允许的重试次数为该值减一~\cite{spark413-config}。

这些参数决定计算尝试何时更换。
某个 attempt 发送 HTTP 后失去响应，后继 attempt 可能再次发送。
代码若用 attempt 编号建立付款键，调度恢复就会形成另一个平台操作。
输入中原 $C$ 与平台 \code{scope} 才定位业务责任，
attempt 编号用于解释由哪个进程执行过哪些步骤。

\code{OutputCommitCoordinator} 在 driver 中保存
stage、stage attempt、分区与获准提交的 task attempt。
executor 请求输出提交许可时，
已经授权的尝试排斥同一分区的其他尝试。
若获准尝试随后失败，协调者可以清除相应授权，允许后继尝试继续。
这一许可状态属于计算调度，输出字节的实际保存由后端提交器与文件系统完成~\cite{spark413-commit}。

\CTwentyFigure{ch20-06}{Spark 区分依赖重算、提交许可与输出后端的保存}{fig:ch20-spark-commit}

只有接入相应文件输出路径的动作经过这个协调点。
普通 \code{foreach} 任务直接写数据库、调用支付或发消息时，
这些外部接口各自接收请求并保存效果。
数据库以业务唯一键和事务关联重复写入，
支付按原 $C/T$ 查询和恢复，
事件以原 $O/EV$ 及消费事务关联。
运行时重算负责再取得计算结果，外部保存域负责识别重复业务动作。

\subsection{发布差异与执行修复}

对账作业先完整发布差异数据，再由独立恢复服务消费。
每条差异包含来源版本、原 $S/C/T$、观察到的平台事实和待核验原因。
恢复者读取当前权威与平台原操作，
再决定补记结果、继续查询或转入人工核对。
旧快照提供发现线索，在线权威提供执行时的资格与当前责任。

例如，$J_1$ 发现 $C_1$ 在平台成功为 $X_1$，
权威导出尚未保存这份结果。
恢复者按原 \code{scope/pay-C1} 查询到 $X_1$，
把完整结果写回原 $C_1/T_1$。
另一个恢复者若已完成回填，重复处理就取得原结果。
$J_2$ 稍后又交付同一条差异时，也通过原 $C_1$ 取得这份结果。
作业分别保存自己的观察，付款记录保存同一次收款的处置。

\CTwentyFigure{ch20-07}{差异文件可以重复交付，业务修复沿原 $C/T$ 汇合}{fig:ch20-repair}

批量结果的保留和回收跟随这些引用。
作业失败留下的尝试文件，在确认没有正式清单引用后可以清理。
已发布清单、输入版本与差异处置记录则保存审计和恢复依据。
修复完成后，原对账结果继续表示当时固定材料中的差异，
新的处置状态另行保存，读者据此区分发现时的观察与后来完成的回填。
